\documentclass[border=10pt,tikz]{standalone}
\usepackage{fontspec}
\setmainfont{Apple SD Gothic Neo}
\setsansfont{Apple SD Gothic Neo}
\usepackage{tikz}
\usetikzlibrary{positioning,arrows.meta,shapes,fit,backgrounds,calc}
\begin{document}
% _design.tex — shared TikZ design system for all blog diagrams
% Inputs nothing on its own — meant to be \input{} from each diagram .tex
% AFTER \begin{document}, BEFORE \begin{tikzpicture}.
%
% Provides a unified visual language across GoF / DSA / EMC++ / Beautiful series:
%   - Colors: muted, accessible, color-coded by role
%   - Node styles: consistent sizing, rounded, subtle borders
%   - Edge styles: UML-style inheritance, aggregation, plus dashed alternates
%
% Usage:
%   \documentclass[border=10pt,tikz]{standalone}
%   \usepackage{tikz}
%   \usetikzlibrary{positioning,arrows.meta,shapes,fit,backgrounds}
%   \begin{document}
%   \input{../_design.tex}     % adjust relative path
%   \begin{tikzpicture}[blog]  % apply 'blog' style
%   ...

\definecolor{absfill}{RGB}{220,232,247}   % cool blue
\definecolor{absbord}{RGB}{72,118,182}
\definecolor{confill}{RGB}{249,222,222}   % warm red
\definecolor{conbord}{RGB}{197,93,93}
\definecolor{impfill}{RGB}{223,238,222}   % soft green
\definecolor{impbord}{RGB}{99,154,93}
\definecolor{clifill}{RGB}{251,239,205}   % pale yellow
\definecolor{clibord}{RGB}{200,167,72}
\definecolor{hifill} {RGB}{253,224,200}   % accent orange
\definecolor{hibord} {RGB}{210,135,68}
\definecolor{nodefill}{RGB}{250,250,250}
\definecolor{nodebord}{RGB}{120,120,120}
\definecolor{edgec}  {RGB}{80,80,80}
\definecolor{edgesc} {RGB}{145,145,145}
\definecolor{groupbord}{RGB}{200,200,200}

\tikzset{
  blog/.style={
    font=\sffamily\small,
    every node/.style={font=\sffamily\small}
  },
  % --- node roles ---
  cls/.style={
    rectangle, rounded corners=2pt,
    draw=nodebord, line width=0.4pt,
    minimum height=8mm, minimum width=28mm,
    inner sep=3pt, fill=nodefill, align=center
  },
  abs/.style={cls, draw=absbord, fill=absfill, font=\sffamily\itshape\small},
  con/.style={cls, draw=conbord, fill=confill},
  imp/.style={cls, draw=impbord, fill=impfill},
  cli/.style={cls, draw=clibord, fill=clifill},
  hi/.style={cls, draw=hibord, fill=hifill, font=\sffamily\bfseries\small},
  % --- circle (for trees / graph nodes) ---
  vx/.style={
    circle, draw=absbord, line width=0.4pt,
    minimum size=8mm, fill=absfill, inner sep=0pt
  },
  vxc/.style={vx, draw=conbord, fill=confill},
  vxh/.style={vx, draw=hibord, fill=hifill, font=\sffamily\bfseries},
  % --- decision diamond ---
  q/.style={
    diamond, draw=clibord, line width=0.4pt, fill=clifill,
    inner sep=2pt, aspect=2.4, align=center
  },
  % --- edges ---
  inh/.style={                                  % UML inheritance
    -{Triangle[length=3mm,width=3mm,open]},
    draw=edgec, line width=0.5pt
  },
  agg/.style={                                  % UML aggregation
    {Diamond[length=2.5mm,width=2.5mm,open]}-{Stealth[length=2.2mm]},
    draw=edgec, line width=0.45pt
  },
  uses/.style={                                 % directed reference
    -{Stealth[length=2.2mm]}, draw=edgec, line width=0.5pt
  },
  arr/.style={                                  % plain arrow (graphs/trees)
    -{Stealth[length=2mm]}, draw=edgec, line width=0.45pt
  },
  edge/.style={                                 % undirected (trees)
    draw=edgec, line width=0.45pt
  },
  alt/.style={                                  % dashed: similar/alt
    -{Stealth[length=2mm]}, dashed, draw=edgesc, line width=0.45pt
  },
  ret/.style={                                  % dashed creates/returns
    -{Stealth[length=2mm]}, dashed, draw=edgesc, line width=0.45pt
  },
  thr/.style={                                  % threading link (red dashed)
    ->, dashed, draw=conbord, line width=0.5pt
  },
  % --- group / region ---
  group/.style={
    rectangle, rounded corners=4pt,
    draw=groupbord, line width=0.4pt,
    inner sep=10pt, fill=black!2
  },
  glabel/.style={font=\sffamily\bfseries\small, anchor=north west, inner sep=2pt}
}

% Korean font convenience macro — included by every Hangul diagram
\providecommand{\KOR}{}

% Sources (verified):
%   Linux Documentation/arch/arm64/booting.rst  spin-table: cpu-release-addr, poll (optional wfe),
%       primary issues sev, jump to non-zero value; psci: kernel issues CPU_ON
%   Linux arch/arm64/kernel/smp_spin_table.c   writeq_relaxed(release_addr) + sev()
%   Linux arch/arm64/kernel/psci.c             cpu_psci_cpu_boot(): psci_ops.cpu_on(...)
%   Linux arch/arm64/kernel/smp.c              __cpu_up(): boot_secondary(), wait, cpu_online()
\begin{tikzpicture}[blog,
  bx/.style={cls, minimum width=62mm, minimum height=11mm, align=center, font=\sffamily\small, inner sep=4pt},
  side/.style={font=\sffamily\footnotesize, text=mutec, align=left, anchor=west},
  brl/.style={font=\sffamily\footnotesize, text=inkc, fill=white, inner sep=2pt}
]

\node[cli, bx] (dt) at (0,0) {부팅에 전달된 DT 또는 ACPI에서\\[2pt]CPU별 enable 방식 확인};

% PSCI branch
\node[prp, bx] (p1) at (-4.2,-2.4) {Linux $\to$ firmware: \texttt{CPU\_ON} 요청\\[2pt](SMC 또는 HVC)};
\node[prp, bx] (p2) at (-4.2,-4.2) {firmware·플랫폼 코드가\\[2pt]대상 CPU 전원·reset 처리};

% spin-table branch
\node[con, bx] (s1) at (4.2,-2.4) {Linux가 \texttt{cpu-release-addr}에\\[2pt]entry 주소 기록 후 SEV};
\node[con, bx] (s2) at (4.2,-4.2) {커널 밖에서 대기하던 CPU가\\[2pt]non-zero 값을 읽고 그 주소로 이동};

\node[abs, bx] (en) at (0,-6.2) {secondary CPU가 커널 entry에 진입};
\node[ok,  bx] (on) at (0,-8.0) {Linux가 CPU online 여부 확인};

\draw[flow] (dt.south) -- ++(0,-0.6) -| node[brl, pos=0.75]{\texttt{"psci"}} (p1.north);
\draw[flow] (dt.south) -- ++(0,-0.6) -| node[brl, pos=0.75]{\texttt{"spin-table"}} (s1.north);
\draw[flow] (p1) -- (p2);
\draw[flow] (s1) -- (s2);
\draw[flow] (p2.south) |- ($(en.north)+(-1.2,0.5)$) -- ($(en.north)+(-1.2,0)$);
\draw[flow] (s2.south) |- ($(en.north)+(1.2,0.5)$) -- ($(en.north)+(1.2,0)$);
\draw[flow] (en) -- (on);

% verification points from the post
\node[side] at (-10.0,-2.4) {확인 1\\[2pt]요청 반환 상태\\[2pt](커널·firmware 로그)};
\node[side] at (3.5,-6.2) {확인 2  secondary 진입};
\node[side] at (3.5,-8.0) {확인 3  online 상태 (sysfs·커널 로그)};

\end{tikzpicture}
\end{document}
